\section{Sign, Visibility, and the Logic of Promise}

The stars function first as a sign of multitude. Their numerical excess, not
their luminosity, supplies the explicit analogy: ``so shall thy seed be.''
The conditional clause ``if thou canst number them'' makes human incapacity
part of the sign. Even a cloudless midnight would leave the task unfinished.

Visibility and countability must therefore not be collapsed. At night Abram
could see more stars than he could count. By day he could know that stars are
there without ordinarily distinguishing them. In vision he might be shown
either ordinary or extraordinary heavens. These modalities produce different
spiritual applications, while the stated comparison---a posterity beyond
practical enumeration---remains stable.

The narrative then reports belief. That order joins sensible sign and divine
word without making faith a substitute for reason. Abram is not praised for
believing that darkness is light or that absent evidence is present evidence.
He entrusts himself to the speaker whose promise exceeds the presently
available human prospect. A daylight application can illuminate this structure
only if it retains the distinction between what the text states, what its
sequence suggests, and what a later reader analogically draws from it.

